% !TEX program = pdflatex
% This file is a visual guide to the available box commands.
\documentclass[a4paper,11pt]{article}
\usepackage{../cleannotes}
\documentsetup{
  title       = {CleanNotes},
  subtitle    = {Visual styles and quick reference},
  shorttitle  = {CleanNotes - quick reference},
  author      = {First name Last name}
}

% Example-only helper and custom variant. Not required by the template.
\newcommand{\cmd}[1]{{\small\ttfamily\detokenize{#1}}}
\newaccentbox{custombox}{leafgreen}{Lorem ipsum.}

\begin{document}
\makedocumenthead
Choose a box by its appearance. Titles, colours and content are yours to choose.
Every box uses the same three optional arguments:
\begin{quote}\small
\cmd{\begin{accentbox}[Title][docblue][extra options]}\par
Content \ldots\quad\cmd{\end{accentbox}}
\end{quote}
Omit the options for the defaults. Use \cmd{[]} for an empty title.

\section*{The six visual styles}
\begin{accentbox}[Lorem ipsum.]
Lorem ipsum dolor sit amet, consectetur adipiscing elit. Sed do eiusmod tempor
incididunt ut labore et dolore magna aliqua.
\end{accentbox}
{\small\cmd{accentbox} --- Soft background, coloured left border, inline title.}

\begin{framedbox}[Lorem ipsum]
Lorem ipsum dolor sit amet, consectetur adipiscing elit.
\[a^2+b^2=c^2\]
\end{framedbox}
{\small\cmd{framedbox} --- Thin outline, separate tinted title bar.}

\begin{numberedbox}[Lorem ipsum]\label{box:lorem}
Lorem ipsum dolor sit amet, consectetur adipiscing elit.
\end{numberedbox}
{\small\cmd{numberedbox} --- Automatic number; supports \cmd{\label} and \cmd{\ref}.}

\begin{cardbox}[Lorem ipsum]
Lorem ipsum dolor sit amet, consectetur adipiscing elit. Ut enim ad minim veniam.
\end{cardbox}
{\small\cmd{cardbox} --- Smaller type and padding, title on its own line.}

\begin{solidbox}[Lorem ipsum]
Lorem ipsum dolor sit amet, consectetur adipiscing elit.
\end{solidbox}
{\small\cmd{solidbox} --- Solid colour with white text.}

\begin{splitbox}[Lorem ipsum]\rightheading{Dolor sit amet}
Lorem ipsum dolor sit amet, consectetur adipiscing elit.
\end{splitbox}
{\small\cmd{splitbox} --- Two heading fields; start the body with \cmd{\rightheading{Right text}}.}

\clearpage
\section*{Colours and custom variants}
\begin{minipage}[t]{0.48\linewidth}
\begin{accentbox}[Lorem ipsum.][alertred]
Lorem ipsum dolor sit amet.
\end{accentbox}
\cmd{\begin{accentbox}[Title.][alertred]}
\end{minipage}\hfill
\begin{minipage}[t]{0.48\linewidth}
\begin{accentbox}[Lorem ipsum.][leafgreen]
Lorem ipsum dolor sit amet.
\end{accentbox}
\cmd{\begin{accentbox}[Title.][leafgreen]}
\end{minipage}

Available colours: \textcolor{docblue}{\cmd{docblue}},
\textcolor{tiporange}{\cmd{tiporange}}, \textcolor{alertred}{\cmd{alertred}},
\textcolor{leafgreen}{\cmd{leafgreen}}, \textcolor{slategray}{\cmd{slategray}}
and \textcolor{softgray}{\cmd{softgray}}. Any named xcolor colour also works.

Define a reusable variant once in the preamble:
\par\cmd{\newaccentbox{custombox}{leafgreen}{Lorem ipsum.}}
\begin{custombox}
Lorem ipsum dolor sit amet, consectetur adipiscing elit.
\end{custombox}
Then use \cmd{\begin{custombox}} \ldots\ \cmd{\end{custombox}}.
It accepts the same three optional arguments as the built-in boxes.
For a one-off adjustment, use the third argument:
\par\cmd{\begin{accentbox}[Title.][docblue][colback=white]}

\section*{Other built-in commands}
{\small\renewcommand{\arraystretch}{1.3}
\begin{tabularx}{\linewidth}{@{}lX@{}}
\toprule
\textbf{Command} & \textbf{Purpose or output} \\
\midrule
\cmd{\documentsetup{title={...},author={...}}} & Title-block and PDF metadata. \\
\cmd{\makedocumenthead} & Compact title block, as on page 1. \\
\cmd{\newaccentbox{name}{colour}{Title.}} & Define an accent-box variant. \\
\cmd{\rightheading{Right text}} & Right heading inside \cmd{splitbox}. \\
\cmd{\begin{steps} ... \end{steps}} & Labelled steps: \cmd{\item[Label] Content.} \\
\cmd{\step{Lorem}} & \step{Lorem} ipsum dolor sit amet. \\
\cmd{\reason{lorem ipsum}} & In maths: \(x=2\quad\reason{lorem ipsum}\). \\
\cmd{\begin{proof} ... \end{proof}} & Standard proof with an ending square. \\
\bottomrule
\end{tabularx}}
\begin{steps}
\item[Lorem ipsum] Dolor sit amet, consectetur adipiscing elit.
\item[Dolor sit amet] Ut enim ad minim veniam.
\end{steps}
\begin{proof}
Lorem ipsum dolor sit amet, consectetur adipiscing elit.
\end{proof}

\section*{Language, layout and standard LaTeX}
English is the default. In your project, change the package line in
\cmd{preamble.tex} to
\cmd{\usepackage[dutch]{cleannotes}} for Dutch labels and hyphenation.
Command names stay the same. Your own titles and text remain yours to translate.

Other options: \cmd{[unnumbered]} for headings without numbers;
\cmd{[onepager]} for landscape A4 with narrow margins and no running header/footer.
Options combine: \cmd{[dutch,onepager]}. Content determines the page count.

Ready to use: \cmd{\section}, \cmd{\subsection}, \cmd{\tableofcontents},
\cmd{itemize}, \cmd{enumerate}, \cmd{equation}, \cmd{align}, \cmd{tabularx},
\cmd{figure}, \cmd{\includegraphics}, TikZ, \cmd{multicols}, \cmd{\href} and
\cmd{\url}. Cross-reference: Box~\ref{box:lorem} via \cmd{\ref{box:lorem}}.
See the preamble for optional plotting and bibliography packages.
\end{document}
